\chapter{Prayer, Pranava and Hymns}
\label{ch:prayer}

\emph{Why does a yoga session begin and end with a prayer, and what does the candidate recite?}
The syllabus asks the candidate to recite the pranava and a few hymns and to know what they mean. This chapter gives the opening and closing chants of the common protocol, the meaning of the pranava, and the way to lead a class in recitation.

\begin{unverified}
This chapter is written from a notes block whose fact check is still pending. Read it as a draft. The chants follow the Common Yoga Protocol; the reading of the pranava follows the official guidebook. A certified teacher must review it before print.
\end{unverified}

\section{Why a session begins with prayer}

A session begins with a prayer, or at least with a prayerful mood. The purpose is practical as much as devotional: a few moments of chanting draw the attention away from the day's business, slow the breath, and set a calm and receptive tone for what follows. A class that walks straight from the street into a posture is still busy in the mind; a class that has chanted together has already settled.

The prayer of a yoga class is a call for harmony and well-being, not an instruction in belief. Keep it inclusive. A teacher who explains the meaning before the chant makes it something the class understands, not only something it repeats.

\section{The opening hymn}

The common protocol opens with a verse from the Rigveda:

\begin{quote}
\emph{saṃ gacchadhvaṃ saṃ vadadhvaṃ saṃ vo manāṃsi jānatām;\\
devā bhāgaṃ yathā pūrve saṃjānānā upāsate.}
\end{quote}

In plain words it is a wish that the people present move together, speak together and think in agreement, as the gods of old shared their portion in harmony. It suits a class well: a group that is about to practise together asks first to be of one mind.

\section{The closing: sankalpa and peace chant}

The session ends in two steps. First comes a short sankalpa, a personal resolve: to become a healthy, peaceful, joyful and loving person, and to treat all life as kin. Then comes the peace chant:

\begin{quote}
\emph{sarve bhavantu sukhinaḥ, sarve santu nirāmayāḥ;\\
sarve bhadrāṇi paśyantu, mā kaścid duḥkhabhāg bhavet;\\
om śāntiḥ śāntiḥ śāntiḥ.}
\end{quote}

It means: may all be happy, may all be free from illness, may all see what is auspicious, and may no one suffer. The class leaves with a wish for others, not only for itself.

\section{Pranava: the syllable om}

The pranava is the syllable om, also written AUM. The official guidebook, following Patanjali, calls it the word that designates Ishvara. Its japa, repeated with full attention and with its meaning in mind, is said to remove obstacles and turn the mind inward. This is the tradition's own claim; present it as such.

A popular traditional reading splits the syllable into three sounds: A for creation, U for preservation and M for dissolution. The guidebook does not give this reading, so teach it as a common explanation and not as a definition.

\paragraph{Mantra japa.} Japa is the repeated chanting of a mantra. In Patanjali's sense the japa of om is done with its meaning in mind. The Board syllabus lists mantra and japa among the practices leading to meditation, so the same skill returns in the meditation chapter (see \cref{ch:meditation}).

\paragraph{Hymns.} A hymn here means a sacred verse used in ritual and devotional practice. Vedic hymns are valued as ancient sacred sound and are chanted, in the tradition's own words, for purification and blessing. Two facts are often asked: the Gayatri mantra is traditionally attributed to the seer Vishwamitra, and the Upanishadic phrase \emph{aham brahmāsmi} means ``I am Brahman'', the ultimate reality.

\section{Leading recitation in a class}

Say the words slowly, a short phrase at a time, and let the class repeat each phrase. Give the meaning in the language of the class first, so that the chant is understood. Keep the pitch comfortable for the whole group, and let the sound end in a short silence before the next step of the session.

\ycbfig{ycb1-fig-hymns}
\ycbfig{ycb1-fig-pranav}

\begin{recall}
\begin{itemize}
\item Open with a prayer or prayerful mood, to calm the mind and set the tone.
\item Opening hymn: \emph{saṃ gacchadhvaṃ}, a wish to move, speak and think together.
\item Close with a sankalpa, then the peace chant \emph{sarve bhavantu sukhinaḥ}.
\item Pranava: om, the word for Ishvara in the Yoga Sutra; japa with its meaning in mind.
\item A, U, M as creation, preservation and dissolution is a popular reading, not in the guidebook.
\end{itemize}
\end{recall}

\subsection*{Worked examples}
\begin{itemize}
\item \textbf{A student says the prayer feels religious and asks to skip it.} Explain that the opening hymn is a wish for the group to be in harmony and the closing chant a wish for everyone's well-being. Invite the student to sit quietly and share the mood if they prefer not to chant. Nobody is required to chant.
\item \textbf{Model answer: the meaning of the peace chant.} May all be happy; may all be free from illness; may all see what is auspicious; may no one suffer. It closes the session with a wish for others.
\end{itemize}

\subsection*{Check yourself}
\begin{enumerate}
\item What opens a yoga session in the common protocol?
\item What two steps close the session?
\item What does the pranava designate in the Yoga Sutra?
\item What is mantra japa?
\item What are the three parts of AUM in the traditional reading, and where does that reading come from?
\item To which seer is the Gayatri mantra traditionally attributed?
\item What does \emph{aham brahmāsmi} mean?
\end{enumerate}

\subsection*{Answers}
\begin{enumerate}
\item A prayer or prayerful mood, with the harmony hymn \emph{saṃ gacchadhvaṃ}.
\item A sankalpa (personal resolve), then the peace chant.
\item Ishvara.
\item The repeated chanting of a mantra, here om, with its meaning in mind.
\item A (creation), U (preservation), M (dissolution); it is a popular traditional reading and is not in the guidebook.
\item Vishwamitra.
\item ``I am Brahman'', the ultimate reality.
\end{enumerate}

\paragraph{Sources.} Common Yoga Protocol, Ministry of AYUSH (opening hymn, sankalpa, peace chant); the official guidebook for Level I (2016, older scheme) on Ishvara pranidhana and japa; the Board syllabus for Level 1.
